\ifdefined\MyWebBook\else
\documentclass[../textbook.tex]{subfiles}
\begin{document}
\fi
\bookchapter{扩散模型的概率原理}{ch:diffusion-probability}

自回归模型沿序列位置逐步生成，扩散模型沿噪声尺度逐步生成。\term{扩散模型}{Diffusion Model}{}先规定一个容易采样的扰动过程，再学习与其对应的逆向生成过程。扩散模型以潜变量概率模型为基础：观测数据被连接到一系列带噪潜变量，训练通过可计算的条件分布约束反向转移。

本章以\term{去噪扩散概率模型}{Denoising Diffusion Probabilistic Model}{DDPM}为主要构造\cite{ho2020ddpm}，逐步推导前向边缘、条件后验和训练目标，并讨论不同输出参数化与采样路径。\bookxref{ch:conditional-vision}再讨论文本条件、潜空间、网络主干和视觉生成系统，分别说明概率过程、网络结构与应用条件。

\section{前向扩散过程}

\subsection{符号与假设}
设数据 $\vect{x}_0\in\mathbb R^D$ 服从数据分布 $q_{\mathrm{data}}$。图像可以按通道和空间位置展平为 $D$ 维向量；所有推导逐元素成立，不要求网络实际展平图像。带噪变量 $\vect{x}_1,\ldots,\vect{x}_T$ 与数据同形，$t\in\{1,\ldots,T\}$ 表示离散噪声步，$t=0$ 专指干净端点。

\begin{table}[htbp]
\centering
\caption{扩散概率模型主要符号}\label{tab:diffusion-symbols}
\begin{tabular}{ll}
\toprule 符号 & 含义\\
\midrule
$D,T$ & 单个样本的标量维数和前向链总步数。\\
$\beta_t,\alpha_t$ & 每步噪声方差，以及信号方差保留系数 $\alpha_t=1-\beta_t$。\\
$\bar\alpha_t$ & 累计保留系数 $\prod_{s=1}^t\alpha_s$，约定 $\bar\alpha_0=1$。\\
$a_t,b_t$ & 边缘信号和噪声标准差，$a_t=\sqrt{\bar\alpha_t}$、$b_t=\sqrt{1-\bar\alpha_t}$。\\
$\vect{\epsilon},\vect{\xi}_t$ & 标准高斯随机变量；边缘总噪声与单步噪声应区分。\\
$\widetilde{\vect{\mu}}_t,\widetilde\beta_t$ & 给定干净样本时反向条件后验的均值和标量方差。\\
$\vect{\mu}_{\vect{\theta}},\sigma_t^2$ & 模型反向条件分布的均值及固定方差。\\
$\vect{\epsilon}_{\vect{\theta}},\widehat{\vect{x}}_0,\vect{v}_{\vect{\theta}}$ & 噪声预测、干净样本预测和速度参数化预测。\\
\bottomrule
\end{tabular}
\end{table}

\begin{assumption}[离散高斯扩散的推导条件]
本章取固定日程 $0<\beta_t<1$，前向各步使用相互独立且与 $\vect{x}_0$ 独立的标准高斯噪声。除输出端点另作说明，反向模型取具有正方差的各向同性高斯分布。数据预处理和尺度保持固定。高斯条件后验公式针对 $t\ge2$；$t=1$ 给定 $\vect{x}_0$ 后发生退化，非退化高斯 KL 公式在那里失去适用前提。
\end{assumption}

\subsection{递推扰动}
\term{前向过程}{Forward Process}{}定义为
\begin{equation}\label{eq:diff-forward}
 \begin{aligned}
  q(\vect{x}_{1:T}\mid \vect{x}_0)&=\prod_{t=1}^{\mathrm{T}}q(\vect{x}_t\mid \vect{x}_{t-1}),\\
  q(\vect{x}_t\mid \vect{x}_{t-1})&=\mathcal N(\sqrt{\alpha_t}\vect{x}_{t-1},\beta_t \mat{I}_D).
 \end{aligned}
\end{equation}
等价的采样表示为 $\vect{x}_t=\sqrt{\alpha_t}\vect{x}_{t-1}+\sqrt{\beta_t}\vect{\xi}_t$，其中 $\vect{\xi}_t\sim\mathcal N(0,\mat{I}_D)$。前一状态被缩放后再叠加噪声，使单位方差高斯保持单位方差；一般数据的方差则随累计日程变化。

若原数据均值与协方差分别为 $\vect{\mu}_0,\mat{\Sigma}_0$，则后面将得到
\begin{equation}
 \begin{aligned}
  \mathbb E[\vect{x}_t]&=a_t\vect{\mu}_0,\\
  \operatorname{Cov}(\vect{x}_t)&=\bar\alpha_t\mat{\Sigma}_0+(1-\bar\alpha_t)\mat{I}_D.
 \end{aligned}
\end{equation}
因此噪声化会逐渐削弱原数据的均值、尺度和相关结构。只有累计信号足够小，末端分布才适合用标准高斯近似；步数多少与这项尺度判断没有直接关系。

\begin{figure}[htbp]
\centering
\begin{tikzpicture}[>=stealth,node distance=14mm,every node/.style={font=\small},b/.style={draw,rounded corners,minimum width=18mm,minimum height=10mm,align=center}]
\node[b] (x0){$\vect{x}_0$\\数据};\node[b,right=of x0] (x1){$\vect{x}_1$};\node[b,right=of x1] (xm){$\cdots$};\node[b,right=of xm] (xt){$\vect{x}_{T-1}$};\node[b,right=of xt] (noise){$\vect{x}_T$\\近似噪声};
\draw[->] (x0) to[bend left=22] node[above]{$q(\vect{x}_1\mid\vect{x}_0)$} (x1);\draw[->] (x1) to[bend left=22] node[above]{$q$} (xm);\draw[->] (xm) to[bend left=22] node[above]{$q$} (xt);\draw[->] (xt) to[bend left=22] node[above]{$q(\vect{x}_T\mid\vect{x}_{T-1})$} (noise);
\draw[->] (noise) to[bend left=22] node[below]{$p_{\vect{\theta}}(\vect{x}_{T-1}\mid\vect{x}_T)$} (xt);\draw[->] (xt) to[bend left=22] node[below]{$p_{\vect{\theta}}$} (xm);\draw[->] (xm) to[bend left=22] node[below]{$p_{\vect{\theta}}$} (x1);\draw[->] (x1) to[bend left=22] node[below]{$p_{\vect{\theta}}(\vect{x}_0\mid\vect{x}_1)$} (x0);
\end{tikzpicture}
\caption[扩散模型正向加噪与反向去噪过程]{扩散模型正向加噪与反向去噪过程。正向加噪马尔可夫链 $q$ 逐步注入高斯噪声，与反向参数化去噪生成链 $p_{\vect{\theta}}$ 逐步去噪生成。}\label{fig:diffusion-markov-chain}
\end{figure}

\section{前向边缘分布的闭式}

\subsection{累积噪声}
将前向递推连续代入，得到
\begin{equation}
 \vect{x}_t=\left(\prod_{r=1}^t\sqrt{\alpha_r}\right)\vect{x}_0+
 \sum_{s=1}^t\left(\sqrt{\beta_s}\prod_{r=s+1}^t\sqrt{\alpha_r}\right)\vect{\xi}_s.
\end{equation}
空乘积取1。独立高斯的线性组合仍为高斯，因此条件均值是 $\sqrt{\bar\alpha_t}\vect{x}_0$，条件协方差是标量 $V_t$ 乘 $\mat{I}_D$，其中
\begin{equation}
 V_t=\sum_{s=1}^t\beta_s\prod_{r=s+1}^t\alpha_r.
\end{equation}
用 $\beta_s=1-\alpha_s$ 将每项写成两个乘积之差，求和发生望远镜消去，得到 $V_t=1-\prod_{s=1}^t\alpha_s=1-\bar\alpha_t$。也可以使用递推 $V_t=\alpha_tV_{t-1}+\beta_t$、$V_0=0$ 归纳证明。

于是\term{前向边缘分布}{Forward Marginal Distribution}{}为
\begin{equation}\label{eq:diff-marginal}
 \begin{aligned}
  q(\vect{x}_t\mid \vect{x}_0)&=\mathcal N(a_t\vect{x}_0,b_t^2\mat{I}_D),\\
  \vect{x}_t&=a_t\vect{x}_0+b_t\vect{\epsilon},\quad\vect{\epsilon}\sim\mathcal N(0,\mat{I}_D).
 \end{aligned}
\end{equation}
式\eqref{eq:diff-marginal}中的 $\vect{\epsilon}$ 为累计高斯扰动的归一化标准差变量，不同于式\eqref{eq:diff-forward}单步转移中的增量噪声 $\vect{\xi}_t$。训练仅需采样特定时间步的边缘配对 $(\vect{x}_0,\vect{x}_t)$，无需顺序迭代生成中间潜变量 $\vect{x}_1,\ldots,\vect{x}_{t-1}$。

\subsection{同边缘分布的路径差异}
固定单一样本配对 $(\vect{x}_0,\vect{\epsilon})$ 并随时间 $t$ 绘制 $a_t\vect{x}_0+b_t\vect{\epsilon}$，可直观呈现信噪比演变趋势；该轨迹表示共享同一噪声基底的确定性插值，并不对应式\eqref{eq:diff-forward}中带独立布朗增量的高斯链样本路径。边缘分布的逐时刻一致性并不等价于联合概率路径的同一性。

这一差别对采样器尤其重要：训练可以只依赖边缘加噪配对，而不同反向生成路径可能共享这些训练配对。后面的 DDIM 正是利用这种区分构造不同的采样过程。

\subsection{信噪比和末端先验}
以数据单位方差为比较口径，\term{信噪比}{Signal-to-Noise Ratio}{SNR}为
\begin{equation}
 \operatorname{SNR}_t=\frac{a_t^2}{b_t^2}=\frac{\bar\alpha_t}{1-\bar\alpha_t}.
\end{equation}
当原数据方差不等于1时，某坐标的实际信号噪声方差比还要乘该坐标的数据方差。因此预处理尺度会改变相同日程对应的学习任务。

取末端先验 $p(\vect{x}_T)=\mathcal N(0,\mat{I}_D)$，条件末端 KL 可以直接计算：
\begin{equation}
 D_{\mathrm{KL}}(q(\vect{x}_T\mid \vect{x}_0)\|p(\vect{x}_T))=
 \frac12\left[\bar\alpha_T\|\vect{x}_0\|^2+D\{-\bar\alpha_T-\log(1-\bar\alpha_T)\}\right].
 \label{eq:diff-terminalkl}
\end{equation}
固定 $D$ 且数据具有有限二阶矩时，$\bar\alpha_T\to0$ 使其期望趋于零。有限 $T$ 下直接从标准高斯开始采样仍有先验近似误差；数据维数和尺度也会影响该误差，而不仅是单个系数看起来是否接近零。

\section{反向条件后验的配方推导}

\subsection{可计算的后验与未知逆分布}
生成希望得到 $q(\vect{x}_{t-1}\mid \vect{x}_t)$，但它需要对未知的干净数据后验积分：
\begin{equation}
 q(\vect{x}_{t-1}\mid \vect{x}_t)=\int q(\vect{x}_{t-1}\mid \vect{x}_t,\vect{x}_0)q(\vect{x}_0\mid \vect{x}_t)\,\mathrm d\vect{x}_0.
\end{equation}
即使给定 $\vect{x}_0$ 的条件是高斯，该混合分布通常偏离单一高斯。模型采用高斯反向转移作为近似族；逆转移的具体形状由数据分布与混合权重共同决定。

训练时 $\vect{x}_0$ 已知，因此可使用\term{反向条件后验}{Reverse Conditional Posterior}{} $q(\vect{x}_{t-1}\mid \vect{x}_t,\vect{x}_0)$。令 $\vect{y}=\vect{x}_{t-1}$，根据前向马尔可夫性质和贝叶斯公式，
\begin{equation}
 q(\vect{y}\mid \vect{x}_t,\vect{x}_0)\propto q(\vect{x}_t\mid \vect{y})q(\vect{y}\mid \vect{x}_0).
\end{equation}
取负对数并去掉与 $\vect{y}$ 无关的常数，得到
\begin{equation}
 -2\log q(\vect{y}\mid \vect{x}_t,\vect{x}_0)=\frac{\|\vect{x}_t-\sqrt{\alpha_t}\vect{y}\|^2}{\beta_t}
 +\frac{\|\vect{y}-\sqrt{\bar\alpha_{t-1}}\vect{x}_0\|^2}{1-\bar\alpha_{t-1}}+C.
\end{equation}
展开平方，整理为 $A_t\|\vect{y}\|^2-2\vect{h}_t^{\mathrm{T}}\vect{y}+C'$，其中
\begin{align}
 A_t&=\frac{\alpha_t}{\beta_t}+\frac1{1-\bar\alpha_{t-1}}
 =\frac{1-\bar\alpha_t}{\beta_t(1-\bar\alpha_{t-1})},\\
 \vect{h}_t&=\frac{\sqrt{\alpha_t}}{\beta_t}\vect{x}_t+
 \frac{\sqrt{\bar\alpha_{t-1}}}{1-\bar\alpha_{t-1}}\vect{x}_0.
\end{align}
$A_t$ 是标量精度，$\vect{h}_t\in\mathbb R^D$。配方为
\begin{equation}
 A_t\|\vect{y}-A_t^{-1}\vect{h}_t\|^2+C'',
\end{equation}
因此后验方差与均值分别为
\begin{equation}
 \widetilde\beta_t=A_t^{-1}=\frac{1-\bar\alpha_{t-1}}{1-\bar\alpha_t}\beta_t.
 \label{eq:diff-postvar}
\end{equation}
\begin{equation}
 \widetilde{\vect{\mu}}_t=\frac{\sqrt{\bar\alpha_{t-1}}\beta_t}{1-\bar\alpha_t}\vect{x}_0+
 \frac{\sqrt{\alpha_t}(1-\bar\alpha_{t-1})}{1-\bar\alpha_t}\vect{x}_t.
 \label{eq:diff-postmean}
\end{equation}
两项系数来自精度加权，与图像插值比例的凑配无关；二者之和一般也不为1。

\subsection{端点退化}
当 $t=1$ 时，给定 $\vect{x}_0$ 则条件变量 $\vect{x}_{t-1}=\vect{x}_0$ 确定退化为以 $\vect{x}_0$ 为支撑的狄拉克测度；式\eqref{eq:diff-postvar}方差趋于零，无法直接套用连续高斯 KL 散度公式。训练中将该端点分离为数据对数似然重建项，采样时则按离散网格取值约定输出。\footnote{针对离散像素数据，严格对数似然需在像素离散量化区间积分连续概率密度，而非直接取单点高斯密度值。本章统一采用连续实数概率测度口径展开推导。}

\section{变分目标}

\subsection{Jensen 不等式与下界}
\term{反向过程}{Reverse Process}{}定义模型联合分布
\begin{equation}
 p_{\vect{\theta}}(\vect{x}_{0:T})=p(\vect{x}_T)\prod_{t=1}^{\mathrm{T}}p_{\vect{\theta}}(\vect{x}_{t-1}\mid \vect{x}_t).
\end{equation}
直接积分得到 $p_{\vect{\theta}}(\vect{x}_0)$ 很困难。引入已知的 $q(\vect{x}_{1:T}\mid \vect{x}_0)$，有
\begin{align}
 \log p_{\vect{\theta}}(\vect{x}_0)&=\log\mathbb E_{q(x_{1:T}\mid x_0)}
 \left[\frac{p_{\vect{\theta}}(\vect{x}_{0:T})}{q(\vect{x}_{1:T}\mid \vect{x}_0)}\right]\\
 &\ge\mathbb E_q[\log p_{\vect{\theta}}(\vect{x}_{0:T})-\log q(\vect{x}_{1:T}\mid \vect{x}_0)].
\end{align}
右侧称为\term{证据下界}{Evidence Lower Bound}{ELBO}。本章令 $\mathcal L_{\mathrm{VLB}}$ 为其负值，因而 $-\log p_{\vect{\theta}}(\vect{x}_0)\le\mathcal L_{\mathrm{VLB}}(\vect{x}_0)$。所谓“下界”作用于对数似然，取负号后是负对数似然的上界。

差距可以写成
\begin{equation}
 \mathcal L_{\mathrm{VLB}}(\vect{x}_0)+\log p_{\vect{\theta}}(\vect{x}_0)
 =D_{\mathrm{KL}}(q(\vect{x}_{1:T}\mid \vect{x}_0)\|p_{\vect{\theta}}(\vect{x}_{1:T}\mid \vect{x}_0))\ge0.
\end{equation}
变分目标同时包含负对数似然与后验 KL 差距。

\subsection{变分分解}
固定 $\vect{x}_0$ 后，前向路径可以反向分解为
\begin{equation}
 q(\vect{x}_{1:T}\mid \vect{x}_0)=q(\vect{x}_T\mid \vect{x}_0)\prod_{t=2}^{\mathrm{T}}q(\vect{x}_{t-1}\mid \vect{x}_t,\vect{x}_0).
\end{equation}
把它代入负下界，按相同变量的对数比组合，得到
\begin{align}
 \mathcal L_{\mathrm{VLB}}(\vect{x}_0)=\;&\underbrace{D_{\mathrm{KL}}(q(\vect{x}_T\mid \vect{x}_0)\|p(\vect{x}_T))}_{L_T}\\
 &+\sum_{t=2}^{\mathrm{T}}\underbrace{\mathbb E_{q(x_t\mid x_0)}D_{\mathrm{KL}}(q(\vect{x}_{t-1}\mid \vect{x}_t,\vect{x}_0)\|p_{\vect{\theta}}(\vect{x}_{t-1}\mid \vect{x}_t))}_{L_{t-1}}\\
 &\underbrace{-\mathbb E_{q(x_1\mid x_0)}\log p_{\vect{\theta}}(\vect{x}_0\mid \vect{x}_1)}_{L_0}.
 \label{eq:diff-elbo}
\end{align}
$L_T$ 将加噪终点与先验对齐；$L_{t-1}$ 约束每一步逆转移；$L_0$ 描述生成端点如何赋予数据概率。固定日程和先验时，$L_T$ 不依赖 $\vect{\theta}$，可以从梯度计算中省去，但计算完整似然界时仍须计入。

\subsection{固定方差下的均值误差}
令 $p_{\vect{\theta}}(\vect{x}_{t-1}\mid \vect{x}_t)=\mathcal N(\vect{\mu}_{\vect{\theta}}(\vect{x}_t,t),\sigma_t^2\mat{I}_D)$，且 $\sigma_t^2>0$ 固定。对 $t\ge2$，高斯 KL 为
\begin{equation}
 D_{\mathrm{KL}}=
 \frac{\|\widetilde{\vect{\mu}}_t-\vect{\mu}_{\vect{\theta}}\|^2}{2\sigma_t^2}
 +\frac D2\left[\frac{\widetilde\beta_t}{\sigma_t^2}-1+\log\frac{\sigma_t^2}{\widetilde\beta_t}\right].
\end{equation}
第二部分在上述假设下与 $\vect{\theta}$ 无关，所以均值学习化为加权平方误差。如果反向方差也由网络预测，比例项与对数项一般都依赖参数，删除它们之后优化目标已不等于完整变分界。

\section{去噪模型的参数化}

\subsection{噪声预测目标}
由 $\vect{x}_t=a_t\vect{x}_0+b_t\vect{\epsilon}$ 得 $\vect{x}_0=\frac{\vect{x}_t-b_t\vect{\epsilon}}{a_t}$。将其代入式\eqref{eq:diff-postmean}，化简为
\begin{equation}
 \widetilde{\vect{\mu}}_t=\frac1{\sqrt{\alpha_t}}\left(\vect{x}_t-\frac{\beta_t}{b_t}\vect{\epsilon}\right).
\end{equation}
于是用\term{噪声预测}{Noise Prediction}{}网络 $\vect{\epsilon}_{\vect{\theta}}(\vect{x}_t,t)$ 定义
\begin{equation}
 \vect{\mu}_{\vect{\theta}}(\vect{x}_t,t)=\frac1{\sqrt{\alpha_t}}\left(\vect{x}_t-\frac{\beta_t}{b_t}\vect{\epsilon}_{\vect{\theta}}(\vect{x}_t,t)\right).
 \label{eq:diff-epsmean}
\end{equation}
两者相减并平方，得到依赖参数的逐步项
\begin{equation}\label{eq:diff-weight}
 \begin{aligned}
  L_{t-1}&=\mathbb E\left[w_t\|\vect{\epsilon}-\vect{\epsilon}_{\vect{\theta}}(\vect{x}_t,t)\|^2\right]+C_t,\\
  w_t&=\frac{\beta_t^2}{2\sigma_t^2\alpha_t(1-\bar\alpha_t)}.
 \end{aligned}
\end{equation}
这里期望包含干净数据和加噪变量；$C_t$ 与均值参数无关。噪声预测由概率目标重新参数化得到：网络根据带噪输入和时间步估计加噪变量，同一带噪输入可以对应不同的干净样本与噪声组合。

\subsection{简化损失的权重效应}
常用\term{简化噪声损失}{Simplified Noise Loss}{}为
\begin{equation}
 \mathcal L_{\mathrm{simple}}=\mathbb E_{x_0,\,t\sim\operatorname{Unif}\{1,\ldots,T\},\,\epsilon}
 \|\vect{\epsilon}-\vect{\epsilon}_{\vect{\theta}}(a_t\vect{x}_0+b_t\vect{\epsilon},t)\|^2.
\end{equation}
它省去或修改了式\eqref{eq:diff-weight}中的时间权重，并将端点也纳入噪声回归。完整变分界还包括对应的时间权重与端点概率项。

若希望用非均匀时间分布 $r_t>0$ 估计内部时间项之和，则
\begin{equation}
 \sum_{t=2}^{\mathrm{T}}\mathbb E[\ell_t]=\mathbb E_{t\sim r}\left[\frac{\ell_t}{r_t}\right],
\end{equation}
其中 $r$ 在 $\{2,\ldots,T\}$ 上归一化，$\ell_t=w_t\|\vect{\epsilon}-\vect{\epsilon}_{\vect{\theta}}\|^2$。改变抽样分布却不补偿 $\frac{1}{r_t}$，会\emph{改变目标}，而不仅改变方差。批内再除以 $D$ 是损失尺度约定；报告损失时必须说明是每样本平方和还是每元素均值。

\begin{note}[均方误差目标下的条件期望解与随机细节恢复边界]
给定 $\vect{x}_t,t$，平方损失的总体最优预测为 $\mathbb E[\vect{\epsilon}\mid \vect{x}_t,t]$。因为 $\vect{x}_0$ 未知，同一个带噪输入可能来自多个干净样本，训练过程中某次独立噪声抽样的全部随机细节超出了网络的恢复能力。生成从先验采样，并通过学习的条件转移构造数据样本。
\end{note}

上述结论可由平方展开证明：令 $m=\mathbb E[\vect{\epsilon}\mid \vect{x}_t,t]$，则对任意预测 $f$，条件风险为 $\mathbb E[\|\vect{\epsilon}-m\|^2\mid \vect{x}_t,t]+\|m-f\|^2$，交叉项因条件均值为零消失。

\subsection{干净样本和速度参数化}
\term{干净样本预测}{Clean Sample Prediction}{}直接输出 $\widehat{\vect{x}}_0$，对应噪声估计为 $\frac{\vect{x}_t-a_t\widehat{\vect{x}}_0}{b_t}$。\term{速度参数化}{Velocity Parameterization}{}定义
\begin{equation}
 \vect{v}=a_t\vect{\epsilon}-b_t\vect{x}_0.
\end{equation}
该参数化见渐进蒸馏研究\cite{salimans2022progressive}。因 $a_t^2+b_t^2=1$，有一个正交变换及其逆：
\begin{equation}
 \begin{aligned}
  \begin{bmatrix}\vect{x}_t\\\vect{v}\end{bmatrix}&=
  \begin{bmatrix}a_t&b_t\\-b_t&a_t\end{bmatrix}
  \begin{bmatrix}\vect{x}_0\\\vect{\epsilon}\end{bmatrix},\\
  \widehat{\vect{x}}_0&=a_t\vect{x}_t-b_t\vect{v}_{\vect{\theta}},\\
  \widehat{\vect{\epsilon}}&=b_t\vect{x}_t+a_t\vect{v}_{\vect{\theta}}.
 \end{aligned}
\end{equation}
令 $a=\cos\phi,b=\sin\phi$，则对固定 $(\vect{x}_0,\vect{\epsilon})$，$\frac{\mathrm d\vect{x}}{\mathrm d\phi}=-b\vect{x}_0+a\vect{\epsilon}=\vect{v}$。这里的速度是相对于角度 $\phi$ 的导数。线性流匹配使用另一条路径，其速度为 $\vect{\epsilon}-\vect{x}_0$。

参数化之间在非退化端点可以转换，但无权重平方损失在各参数化下对应不同的目标。对同一固定 $\vect{x}_t$ 的预测误差，
\begin{equation}
 \|\widehat{\vect{\epsilon}}-\vect{\epsilon}\|^2=\frac{a_t^2}{b_t^2}\|\widehat{\vect{x}}_0-\vect{x}_0\|^2,
 \quad
 \|\widehat{\vect{x}}_0-\vect{x}_0\|^2=b_t^2\|\vect{v}_{\vect{\theta}}-\vect{v}\|^2,
 \quad
 \|\widehat{\vect{\epsilon}}-\vect{\epsilon}\|^2=a_t^2\|\vect{v}_{\vect{\theta}}-\vect{v}\|^2.
\end{equation}
因此等权预测 $\vect{x}_0$、$\vect{\epsilon}$ 或 $\vect{v}$ 隐含不同时间权重。靠近 $a_t=0$ 时，从噪声预测除以 $a_t$ 恢复 $\vect{x}_0$ 可能放大误差；靠近 $b_t=0$ 时，从 $\vect{x}_0$ 反推噪声也有相应问题。有限精度、端点选择和裁剪规则均属于采样协议。三种核心参数化目标的数学形式与数值性质对比见表\ref{tab:diffusion-objectives}。

\begin{table}[htbp]
\centering
\small
\setlength{\tabcolsep}{3pt}
\caption{扩散模型不同参数化训练目标对比}\label{tab:diffusion-objectives}
\begin{tabularx}{\linewidth}{@{}p{1.8cm}p{1.6cm}p{2.2cm}p{2.2cm}X@{}}
\toprule
参数化方案 & 网络输出 & 重构 $\widehat{\vect{x}}_0$ & 重构 $\widehat{\vect{\epsilon}}$ & 数值特性与端点稳定性 \\
\midrule
噪声预测 ($\vect{\epsilon}$) & $\vect{\epsilon}_{\vect{\theta}}(\vect{x}_t, t)$ & $\frac{\vect{x}_t - b_t\vect{\epsilon}_{\vect{\theta}}}{a_t}$ & $\vect{\epsilon}_{\vect{\theta}}$ & 偏向高信噪比；$a_t\to0$（重噪声端）重构数据产生除零敏感 \\
数据预测 ($\vect{x}_0$) & $\widehat{\vect{x}}_{\vect{\theta}}(\vect{x}_t, t)$ & $\widehat{\vect{x}}_{\vect{\theta}}$ & $\frac{\vect{x}_t - a_t\widehat{\vect{x}}_{\vect{\theta}}}{b_t}$ & 偏向低信噪比；$b_t\to0$（微噪声端）反推噪声易数值溢出 \\
速度预测 ($v$) & $\vect{v}_{\vect{\theta}}(\vect{x}_t, t)$ & $a_t\vect{x}_t - b_t\vect{v}_{\vect{\theta}}$ & $b_t\vect{x}_t + a_t\vect{v}_{\vect{\theta}}$ & 正交旋转参数化；全时域无除零奇异点，利于快速采样 \\
\bottomrule
\end{tabularx}
\end{table}

\section{扩散模型的训练}

时间条件告诉网络当前噪声尺度。可以把离散 $t$ 映射到正弦和余弦特征，再通过多层感知机注入主干，也可以使用对数信噪比 $\lambda_t=\log\left(\frac{a_t^2}{b_t^2}\right)$。索引相同但日程不同，会对应不同恢复任务；模型输入的时间单位、归一化与日程必须共同确定。端点处对数信噪比可能发散，应采用训练定义的有限范围或专门端点处理。

对图像批量，$\vect{x}_0,\vect{\epsilon},\vect{x}_t$ 通常形状为 $[B,C,H,W]$，时间索引为 $[B]$。取出的 $a_t,b_t$ 重排为 $[B,1,1,1]$ 后广播；预测输出应与 $\vect{x}_t$ 同形。概率推导不限定主干采用卷积 U-Net 还是 Transformer，它们改变函数表达与计算结构，不改变加噪变量的定义。

\begin{figure}[htbp]
\centering
\begin{tikzpicture}[>=stealth,node distance=8mm,every node/.style={font=\small},b/.style={draw,rounded corners,text width=27mm,minimum height=11mm,align=center}]
\node[b] (data){数据 $\vect{x}_0$};
\node[b,right=of data] (noisy){闭式加噪\\$\vect{x}_t=a_t\vect{x}_0+b_t\vect{\epsilon}$};
\node[b,right=of noisy] (net){时间条件网络\\预测 $\vect{\epsilon}_{\vect{\theta}}$};
\node[b,right=of net] (loss){加权误差\\与目标 $\vect{\epsilon}$ 比较};
\node[b,below=of noisy] (random){抽样 $t,\vect{\epsilon}$};
\draw[->] (data)--(noisy);
\draw[->] (noisy)--(net);
\draw[->] (net)--(loss);
\draw[->] (random)--(noisy);
\draw[->] (random)-|(net);
\draw[->] (random)-|(loss);
\end{tikzpicture}
\caption[扩散模型闭式采样与训练闭环]{扩散模型闭式采样与训练闭环。利用重参数化高斯性质闭式采样任意时间步加噪图像，去噪网络预测噪声并反向传播更新。}\label{fig:diffusion-training-loop}
\end{figure}

\begin{algorithm}[噪声预测训练]
\AlgoInput{数据分布、固定日程、预测类型、时间采样分布和损失权重。}
\AlgoOutput{训练后的预测网络。}
\begin{myalgorithmic}[1]
\State 抽取干净样本 $\vect{x}_0$ 并施加固定预处理；独立抽取时间 $t$ 与标准高斯噪声 $\vect{\epsilon}$
\State 直接由累积日程构造 $\vect{x}_t=a_t\vect{x}_0+b_t\vect{\epsilon}$，无需运行完整前向马尔可夫链
\State 输入带指定时间条件的 $\vect{x}_t$；依据预测类型构造目标 $\vect{\epsilon}$、$\vect{x}_0$ 或 $a_t\vect{\epsilon}-b_t\vect{x}_0$
\State 依据选定目标与时间采样规则计算权重，随后反向传播并更新参数；完整变分训练还包含相应端点与方差项
\State 达到预定训练步数或明确定义的停止条件后，输出模型权重，连同日程、预测类型与数据尺度配置
\end{myalgorithmic}
不变量包括：监督目标与输入扰动源于同一噪声实例、时间步嵌入与累计方差日程严格同步、批内样本遵循统一输入归一化尺度。优化迭代步数（Training Steps）与连续扩散时间坐标（Diffusion Timesteps）相互解耦，分别计量参数更新与噪声衰减尺度。
\end{algorithm}

\section{DDIM 采样}

\subsection{随机反向链}
生成不具有已知干净样本。它从 $\vect{x}_T\sim\mathcal N(0,\mat{I}_D)$ 开始，在每一步预测噪声，再依据反向分布更新
\begin{equation}
 \vect{x}_{t-1}=\vect{\mu}_{\vect{\theta}}(\vect{x}_t,t)+\sigma_t\vect{\xi}_t,\qquad \vect{\xi}_t\sim\mathcal N(0,\mat{I}_D).
\end{equation}
反向递推内部步骤的高斯扰动须独立重新采样。在生成终端 $t=1$，工程实现通常直接取后验条件均值输出以降低背景纹理噪声；若严格遵循概率生成规范，则应在终端似然区间内完成一次显式采样，实际部署时应统一声明所采用的截断协议。

选择 $\sigma_t^2=\widetilde\beta_t$ 或其他方差会改变采样过程。只删除部分训练时间索引而继续沿用原来的相邻单步 $\alpha_t,\beta_t$ 更新，系数与时间网格就对不上了。跨步转移需要按两个保留时间之间的累计系数重建。

\subsection{共享边缘的不同路径}
\term{去噪扩散隐式模型}{Denoising Diffusion Implicit Model}{DDIM}构造与相同加噪边缘兼容的另一类路径，并允许确定性采样\cite{song2021ddim}。设从时间 $t$ 跳至更早时间 $s<t$，记模型恢复的干净样本为 $\widehat{\vect{x}}_0=\frac{\vect{x}_t-b_t\widehat{\vect{\epsilon}}}{a_t}$，一种更新写为
\begin{equation}
 \vect{x}_s=a_s\widehat{\vect{x}}_0+\sqrt{b_s^2-\sigma_{t\to s}^2}\,\widehat{\vect{\epsilon}}+\sigma_{t\to s}\vect{\xi},
 \label{eq:diff-ddim}
\end{equation}
其中
\begin{equation}
 \sigma_{t\to s}=\eta\sqrt{\frac{1-\bar\alpha_s}{1-\bar\alpha_t}}
 \sqrt{1-\frac{\bar\alpha_t}{\bar\alpha_s}},\qquad 0\le\eta\le1.
\end{equation}
在此范围及合法日程下，根号项非负。若给定真实 $\vect{x}_0$ 且 $\widehat{\vect{\epsilon}}=\frac{\vect{x}_t-a_t\vect{x}_0}{b_t}$，新状态噪声由两个独立标准高斯的线性组合构成，方差相加为 $b_s^2$，所以保留 $q(\vect{x}_s\mid \vect{x}_0)$ 边缘。这种构造对各时间的联合分布另有选择，与前向马尔可夫链的联合分布可以不同。

$\eta=0$ 时，没有额外随机项，给定初始噪声、网络及时间网格后路径确定。相邻时间且 $\eta=1$ 时，方差成为 $\widetilde\beta_t$，均值与对应后验均值参数化一致。跨步 $\eta=1$ 则使用重构后的粗粒度系数，仅在选定时间点调用模型，其轨迹一般不同于完整原链的逐步采样轨迹。

采样质量由网络误差、离散化路径与时间覆盖共同决定。减少时间步时，应在既定模型上评估采样质量与计算成本。\footnote{确定性采样仍以随机初始噪声为输入，因而输出整体仍可构成分布。确定性指给定初始状态后的映射，不表示所有请求生成同一结果。}

\begin{algorithm}[反向生成]
\AlgoInput{训练后的网络、兼容日程、预测类型、递减时间网格和更新规则。}
\AlgoOutput{生成样本。}
\begin{myalgorithmic}[1]
\State 校验时间网格严格递减，且日程与预测类型兼容；固定推理配置
\State 从指定的终端先验中抽取 $\vect{x}_T$
\For{each 有限时间网格中满足 $t>s\ge0$ 的相邻时间对 $(t,s)$}
  \State 将 $\vect{x}_t,t$ 送入网络，并将其输出转换为更新规则所需的噪声或干净样本估计
  \State 依据 $(t,s)$ 对应的系数计算更新均值与非负方差
  \If{随机采样规则处于激活状态且当前转移需要噪声}
    \State 抽取独立新噪声并依据指定方差更新 $\vect{x}_s$
  \Else
    \State 按确定性或端点规则更新 $\vect{x}_s$
  \EndIf
\EndFor
\State \Return 经逆预处理或下游解码后的生成样本
\end{myalgorithmic}
不变量为时间严格递减、状态与系数对应、预测类型兼容、方差非负。生成路径不读取真实 $x_0$，也不通过生成循环执行参数训练。
\end{algorithm}

\begin{example}

取日程 $\beta_1=0.1,\beta_2=0.2$，则 $\alpha_1=0.9,\alpha_2=0.8$、$\bar\alpha_1=0.9,\bar\alpha_2=0.72$。设 $x_0=1$，边缘噪声 $\epsilon=0.5$，得到
\begin{equation*}
 x_2=\sqrt{0.72}+0.5\sqrt{0.28}\approx1.11310.
\end{equation*}
条件后验精度为 $A_2=\frac{0.8}{0.2}+\frac{1}{0.1}=14$，所以 $\widetilde\beta_2=\frac{1}{14}\approx0.071429$。线性项为
\begin{align*}
 h_2&=\frac{\sqrt{0.8}}{0.2}x_2+\frac{\sqrt{0.9}}{0.1},\\
 \widetilde\mu_2&=\frac{h_2}{14}\approx1.03320.
\end{align*}
亦可由噪声参数化式直接求得 $\widetilde\mu_2=\frac{x_2-\frac{0.2\times0.5}{\sqrt{0.28}}}{\sqrt{0.8}}\approx1.03320$。由于后验方差 $\widetilde\beta_2=\frac{1}{14}$ 保持非零，该单步反向估计输出带噪均值状态，而非直接收敛至确定性真值 $x_0=1$。

若网络预测 $\widehat\epsilon=0.4$，其均值误差为
\begin{equation*}
 \mu_\theta-\widetilde\mu_2=\frac{0.2(0.5-0.4)}{\sqrt{0.8}\sqrt{0.28}}\approx0.04226.
\end{equation*}
选择 $\sigma_2^2=\widetilde\beta_2$，高斯方差项为零，KL 仅为
\begin{equation*}
 \frac{(0.04226)^2}{\frac{2}{14}}\approx0.0125.
\end{equation*}
噪声权重同样给出 $w_2=\frac{0.2^2}{[2(\frac{1}{14})(0.8)(0.28)]}=1.25$，乘噪声误差平方 $(0.1)^2$ 得 $0.0125$。

再以真实噪声估计考察 $2\to1$ 更新。DDIM 在 $\eta=0$ 时给出 $x_1=\sqrt{0.9}+\sqrt{0.1}\times0.5\approx1.10680$；$\eta=1$ 时更新均值约为 $1.03320$，标准差为 $\sqrt{\frac{1}{14}}$。两者采用不同的耦合：前者选择共享总噪声的确定性耦合，后者采用具有条件方差的路径。本例取两步日程，末端累计信号系数 $\bar\alpha_2=0.72$，标准高斯先验与该末端分布仍有较大差距。
\end{example}

\section{扩散模型的误差分析}

前向闭式检查扰动分布，训练损失评价加权回归目标；似然、感知质量与多样性按各自指标评估。反向采样每一步依赖前一步生成状态，模型可能在偏离训练边缘的区域继续运行，误差在那里按相互独立的单步噪声求和会低估实际偏差。

需要保持一致的对象包括数据尺度、累计日程、时间映射、预测类型、反向方差、采样网格和端点规则。预测类型决定输出对应的统计量。对 $\widehat x_0$ 裁剪后，如果更新仍使用未重新关联的噪声估计，也会改变所实现的算法；这种处理必须明确写入采样定义。

前向加噪、时间嵌入、噪声预测和单步均值必须遵循相同的时间步与参数化约定。固定样本与同一总噪声的分镜用于比较各时间尺度；独立增量链的轨迹按前向转移逐步采样。调度器保存日程系数与更新协议。条件 U-Net、文本交叉注意力、无分类器引导和 VAE 属于\bookxref{ch:conditional-vision}的条件生成扩展，流匹配和视频进一步改变路径或数据结构，各扩展使用对应的变量、路径与训练目标。
% BEGIN GENERATED INTERVIEWS

% END GENERATED INTERVIEWS
\chapterreferences
\myendchapterdocument
